\section{Rainbow DQN}
\label{sec:rainbow}

\subsection{Algorithm Overview}

Rainbow~\cite{hessel2018rainbow} combines several independent extensions of the standard DQN into a single, unified algorithm, drastically improving sample efficiency and training stability. We implement all six core Rainbow components: Double Q-learning, Prioritized Experience Replay (PER), Dueling Networks, Multi-step Returns, Distributional Reinforcement Learning (C51), and Noisy Networks. 

Double Q-learning mitigates the overestimation bias inherent to standard Q-learning by decoupling action selection (handled by the online network) from action evaluation (handled by the target network). PER improves data efficiency by prioritizing the sampling of transitions with high temporal-difference errors, ensuring the agent learns from its most surprising experiences. Multi-step returns bootstrap the Q-values over $n$ forward steps, accelerating the propagation of delayed rewards. The dueling architecture splits the neural network into two separate streams (one estimating the state-value and another estimating the state-dependent action advantages) allowing the agent to better evaluate states where specific action choices have minimal impact. Furthermore, distributional RL (C51) models the entire probability distribution of expected returns rather than just their scalar mean, providing a much richer learning signal. Finally, Noisy Networks replace traditional $\epsilon$-greedy exploration by injecting parameterized noise directly into the linear layers of the network, driving state-dependent and mathematically consistent exploration.

\subsection{Architecture and Distributional Calibration}

Applying Rainbow to high-dimensional, pixel-based continuous control requires careful architectural scaling and distributional calibration. In the C51 formulation, the value distribution is modeled over a discrete categorical support bounded by $[v_{\min}, v_{\max}]$. It is critical to properly calibrate this interval to encompass the true range of possible discounted returns. Given our custom reward shaping with a forward weight of $w_{\text{vel}} = 3.0$ and a healthy bonus of $w_{\text{healthy}} = 1.0$, the maximum single-step reward is approximately $r_{\max} \approx 4.0$. 

For an infinite-horizon discounted MDP, the theoretical maximum return is defined as:
\begin{equation}
    R_{\max} \approx \frac{r_{\max}}{1 - \gamma}
\end{equation}
Using our discount factor of $\gamma = 0.99$, we calculate $R_{\max} \approx 400$. Consequently, we calibrated our support bounds to $v_{\min} = -10.0$ and $v_{\max} = 400.0$ distributed across 51 atoms. Poor calibration here severely cripples learning due to distribution clipping. 

Additionally, we found that standard network capacities were insufficient for evaluating pixel-based states in this regime. Increasing the hidden dimension of the dueling head from 256 to 512 neurons proved essential; the smaller configuration failed to accurately capture the complex advantage function, whereas 512 neurons provided the necessary representational capacity. To maintain stability and prevent catastrophic updates, we applied a gradient clipping norm of $10.0$. Note that unlike our standard DQN approach, we do not employ branching architectures here, as extending branching to the Rainbow algorithm yielded no measurable gains for our setup.

% --- FLOATING BOTTOM PLOT FOR HUMANOID ---
\begin{figure*}[t]
    \centering
    \includegraphics[width=\textwidth]{humanoid_rainbow_reward_speed_x.pdf}
    \caption{Learning dynamics of Rainbow on the Humanoid-v5 environment. The three subplots illustrate the steady convergence of reward, velocity, and forward distance.}
    \label{fig:humanoid_rainbow_results}
\end{figure*}
% ---------------------------------------


\subsection{Hopper: A Leap in Performance}

Before tackling the Humanoid, we evaluated our Rainbow configuration on the simpler Hopper environment. The improvements over the baseline BDQN were massive. While standard DQN required approximately 40 million frames to learn a stable gait, Rainbow achieved highly stable walking in just 10 million steps. 



Interestingly, Rainbow also discovered a vastly superior behavioral policy. Previously, the standard DQN converged to a sub-optimal speed of $1.25$ units/s against a target of $1.5$, learning a somewhat fumbling, conservative hop to maintain balance. Rainbow, however, learned to confidently exceed the target, stabilizing at an average velocity of $1.75$ units/s. This behavior suggests that learning a rapid, fluid, and continuous motion is dynamically more stable (relying on forward inertia) than the slower, stuttering hops learned by DQN that constantly require the agent to catch its balance. The Rainbow Hopper sustains this velocity indefinitely, traveling entirely out of the standard rendering zone without collapsing. 

\subsection{Humanoid: Learning Stable Walking Behaviour}

Building on the success in the Hopper environment, we applied Rainbow to the heavily constrained Humanoid-v5 task. The environment utilized a $96 \times 96$ pixel observation with a frame stack of 4, running across 30 parallel environments. The action space was discretized into 3 bins per joint, mapped to the absolute extremes.

\textbf{Algorithm Configuration:}
We trained the model for 50,000,000 frames using a replay buffer of 300,000 transitions and a batch size of 64. We utilized 3-step returns ($n=3$). Prioritized Experience Replay was configured with an $\alpha = 0.6$, and importance sampling weights ($\beta$) were annealed from $0.4$ to $1.0$ over the first 1,000,000 frames. Exploration was driven dynamically via Noisy Networks initialized with $\sigma_0 = 0.5$. The optimization used a learning rate of $10^{-4}$ with target network updates occurring every 10,000 frames. Reward parameters were strictly defined: target velocity at $1.0$, $w_{\text{vel}} = 3.0$, $w_{\text{healthy}} = 1.0$, and a control penalty of $w_{\text{ctrl}} = 0.01$.

\textbf{Results and Discussion:}
Unlike the standard DQN, which completely failed to learn meaningful movement despite exhaustive hyperparameter sweeps, Rainbow successfully achieved stable walking in just 30 million environment steps. Remarkably, it conquered the high-dimensional Humanoid task in fewer steps than the baseline DQN required for the simple Hopper.

As shown in Figure~\ref{fig:humanoid_pose}, the resulting movement is highly stable but visually unnatural. To maximize the reward function while avoiding catastrophic falls, the agent adopts a posture with its torso folded heavily forward. It effectively drags one leg while using the other for primary propulsion, resembling a crutch-like motion. While this motion is physically awkward, it exploits the physics engine brilliantly, keeping the center of mass extremely low and completely bypassing the bipedal instability.